% !TEX root = ./main.tex
% 西财研究生学位论文印刷格式：
% https://graduate.swufe.edu.cn/info/1036/2142.htm
\usepackage{setspace}

% 页眉框顶边距 2.8 cm；页码基线距纸底 2.5 cm。
% 所有版心参数同时交给 geometry 计算，之后不要单独修改 headheight。
\geometry{
  a4paper,
  top=3.9cm, bottom=3.4cm,
  left=3.45cm, right=3.45cm,
  bindingoffset=0pt,
  headheight=14bp,
  headsep=\dimexpr1.1cm-14bp\relax,
  footskip=0.9cm
}

% 22.4 cm 的版心分为 30 个行位；单倍基线经 1.37 倍后落到该网格。
% 小四号汉字宽 12 bp，14.1 cm 的版心每行可容纳 33 个全角汉字。
\newlength{\swufebodybaseline}
\newlength{\swufesinglebaseline}
\newlength{\swufehalflineskip}
\setlength{\swufebodybaseline}{\dimexpr\textheight/30\relax}
\setlength{\swufehalflineskip}{\dimexpr\swufebodybaseline/2\relax}
\ExplSyntaxOn
\dim_set:Nn \swufesinglebaseline
  { \fp_to_dim:n { \dim_to_fp:n { \swufebodybaseline } / 1.37 } }
\ExplSyntaxOff
\AddToHook{cmd/normalsize/after}{%
  \fontsize{12bp}{\swufesinglebaseline}\selectfont
}
\AddToHook{begindocument/end}{%
  \setstretch{1.37}\normalsize\songti
  \setlength{\parskip}{0pt}%
  \setlength{\parindent}{2\ccwd}%
}

% 字体依次使用项目文件、系统同名字体和 TeX Live 替代字体。
% 四种字形齐全才使用项目内的西文字体，避免部分文件缺失时中断编译。
\newif\ifswufeprojecttimes
\swufeprojecttimestrue
\IfFileExists{font/Times New Roman.ttf}{}{\swufeprojecttimesfalse}
\IfFileExists{font/Times New Roman Bold.ttf}{}{\swufeprojecttimesfalse}
\IfFileExists{font/Times New Roman Italic.ttf}{}{\swufeprojecttimesfalse}
\IfFileExists{font/Times New Roman Bold Italic.ttf}{}{\swufeprojecttimesfalse}
\ifswufeprojecttimes
  \setmainfont[
    Path=font/,
    UprightFont={Times New Roman.ttf},
    BoldFont={Times New Roman Bold.ttf},
    ItalicFont={Times New Roman Italic.ttf},
    BoldItalicFont={Times New Roman Bold Italic.ttf}
  ]{}%
\else
  \IfFontExistsTF{Times New Roman}{%
    \setmainfont{Times New Roman}%
  }{%
    % 按文件名交给 TeX 查找，不依赖系统是否注册 TeX Gyre 字体名称。
    \setmainfont[
      Extension=.otf,
      UprightFont=*-regular,
      BoldFont=*-bold,
      ItalicFont=*-italic,
      BoldItalicFont=*-bolditalic
    ]{texgyretermes}%
  }%
\fi

% TeX Live 未提供小标宋；项目文件和系统同名字体均缺失时以 FandolSong 近似。
\IfFileExists{font/FZXiaoBiaoSong.ttf}{%
  \newfontfamily\swufexbslatin[Path=font/]{FZXiaoBiaoSong.ttf}%
  \setCJKfamilyfont{swufexbs}[Path=font/]{FZXiaoBiaoSong.ttf}%
}{%
  \IfFontExistsTF{FZXiaoBiaoSong-B05S}{%
    \newfontfamily\swufexbslatin{FZXiaoBiaoSong-B05S}%
    \setCJKfamilyfont{swufexbs}{FZXiaoBiaoSong-B05S}%
  }{%
    \IfFontExistsTF{方正小标宋简体}{%
      \newfontfamily\swufexbslatin{方正小标宋简体}%
      \setCJKfamilyfont{swufexbs}{方正小标宋简体}%
    }{%
      \newfontfamily\swufexbslatin{FandolSong-Regular.otf}%
      \setCJKfamilyfont{swufexbs}{FandolSong-Regular.otf}%
    }%
  }%
}
\newcommand{\swufexiaobiaosong}{\swufexbslatin\CJKfamily{swufexbs}}

\ctexset{
  chapter={
    name={}, number=\arabic{chapter},
    format=\centering\normalfont\swufexiaobiaosong\zihao{2},
    titleformat=\swufexiaobiaosong\zihao{2},
    numberformat=\swufexiaobiaosong\zihao{2},
    beforeskip=0pt, afterskip=\swufebodybaseline,
    aftername={\ }, pagestyle=fancy
  },
  section={format=\normalfont\heiti\zihao{-3}},
  subsection={format=\normalfont\heiti\zihao{-4}}
}

% 页眉和页码的中文、西文、数字均使用宋体、小五号。
\newfontfamily\swufeheaderlatin{FandolSong-Regular.otf}
\newcommand{\swufeheaderfont}{%
  \normalfont\linespread{1}\swufeheaderlatin\songti\fontsize{9bp}{12bp}\selectfont
}
\makeatletter
\renewcommand{\chaptermark}[1]{\markboth{\arabic{chapter}\ #1}{}}
\fancyhf{}
\fancyhfoffset{0pt}
\setlength{\headwidth}{\textwidth}
\fancyheadinit{\swufeheaderfont}
\fancyfootinit{\swufeheaderfont}
\fancyhead[C]{%
  \swufeheaderfont
  \ifodd\value{page}\leftmark\else\swufezhtitle\fi
}
\fancyfoot[C]{\swufeheaderfont\thepage}
\renewcommand{\headrulewidth}{0.4pt}
\renewcommand{\headrule}{\hrule height 0.4pt width\headwidth}
\let\ps@plain\ps@fancy
\let\ps@headings\ps@fancy
\makeatother
\pagestyle{fancy}

% setspace 会在浮动体开始时恢复 normalsize，因此在环境开始后显式设五号。
\newcommand{\swufefloatfont}{%
  \normalfont\songti\fontsize{10.5bp}{14bp}\selectfont
}
\makeatletter
\apptocmd{\@xfloat}{\swufefloatfont}{}{%
  \PackageError{swufethesis}{Cannot set float font}{Check the setspace version.}%
}
\makeatother
\AddToHook{env/longtable/begin}{\swufefloatfont}
\DeclareCaptionFont{swufefigure}{\swufefloatfont}
\DeclareCaptionFont{swufetable}{%
  \normalfont\heiti\fontsize{10.5bp}{14bp}\selectfont
}
\captionsetup{font=swufefigure,labelfont={},textfont={},skip=6bp}
\captionsetup[table]{font=swufetable,position=top}
\captionsetup*[sub]{font=swufefigure}

% 普通文字与浮动图表之间各留半行；相邻浮动体的两个半行合为一行。
\setlength{\textfloatsep}{\swufehalflineskip}
\setlength{\intextsep}{\swufehalflineskip}
\setlength{\floatsep}{\swufebodybaseline}
\setlength{\dbltextfloatsep}{\swufehalflineskip}
\setlength{\dblfloatsep}{\swufebodybaseline}

\newcommand{\swufetablenotefont}{%
  \normalfont\songti\fontsize{9bp}{12bp}\selectfont
}
\newcommand{\swufetablenote}[1]{%
  \par\smallskip\noindent{\swufetablenotefont #1\par}%
}
\appto\TPTnoteSettings{\swufetablenotefont}

% 参考文献保留小四字号与正文行距，仅去掉列表额外的条目间留白。
\setlength{\bibsep}{0pt}

% 封面不属于正文版心；独立布局避免长标题只将最后一个字挤到下一行。
\newcommand{\swufecovergeometry}{%
  \newgeometry{top=2.8cm,bottom=3cm,left=2.75cm,right=2.75cm,bindingoffset=0pt}%
}
